% #2 MCFG v76 第1段階草稿：局所群型付け，LG の特徴づけ，半束による境界（英語版 appendix_local_group.tex に対応）
\section{局所群型付けと半束による境界}
\label{app:local-group}


\paragraph{半群化の規約（v82移行）．}
この付録の $\mathbf{LG}$ 条件は型付けの
\emph{正語像半群} $S=h(\Sigma^+)$ に課し，
新単位元添加 $S^\sharp$ には課さない．
群準同型とモノイド準同型は計算用の補助表示であり，
意味論上の学習型は正語への制限である．
補題~\ref{lem:semigroup-transfer} により定量的証明と両立する．


付録~\ref{app:window-abelian}，\ref{app:window-group} の窓・群型付けの正像は，
有限局所自明半群と有限群の pseudovariety の join $\mathbf{LI}\vee\mathbf{G}$ に属する~\cite[Lemma~7.4]{kuriyamaCFG}．
ここでは，元文法の厚さによる多項式上界を，それより真に大きい
\emph{局所群}の pseudovariety $\mathbf{LG}$ に拡張する．
$\mathbf{LG}$ は，冪等元 $e$ ごとの局所モノイド $eSe$ がすべて群となる有限半群 $S$ の全体である．
古典的には $\mathbf{LG}=\mathbf G*\mathbf D$ である~\cite{Straubing1985,Tilson1987}．
局所群性，最小イデアルの冪零拡大，および2元半束の不在の同値性は
\emph{既知}の結果である~\cite[Proposition~2.3]{AlmeidaEtAl2009}．
定理~\ref{thm:lg-characterization} はこの同値性を再掲し，
MCFG の定量的評価に必要な窓と局所因子に沿った群積による
型付けの特徴づけを追加する．
特に $\mathbf{LG}$ は，2元半束 $U_1=\{1,\zeta\}$（$\zeta^2=\zeta$）を
含まない有限半群のクラスにちょうど一致する．
この半束は，\cite{kuriyamaCFG} の指数下界を与える型付けの正像である．
したがって pseudovariety の水準では，半束が唯一の障害となる
（系~\ref{cor:pseudovariety-dichotomy}）．

この付録を通じて，$\Sigma$ は空でなく，$k\ge1$，$A$ は有限群とする．
$|w|\ge k+1$ の語 $w$ について，長さ $k+1$ の因子を開始位置の順に並べた列を $\sigma_k(w)$ と書く．

\begin{definition}[局所群型付け]
\label{def:local-group-typing}
$g:\Sigma^{k+1}\to A$ を任意の写像とする．$|w|\ge2k$ に対し，
\[
 \gamma_{k,g}(w):=g(F_1)g(F_2)\cdots g(F_{|w|-k}),
 \qquad \sigma_k(w)=(F_1,\ldots,F_{|w|-k})
\]
を $A$ における順序つき積とし，$|w|<2k$ に対しては $\gamma_{k,g}(w):=1_A$ とする．
\emph{局所群型付け}を
\[
 \lambda_{k,g}(w):=\bigl(h^{k,k}_\Sigma(w),\gamma_{k,g}(w)\bigr)
\]
で定める．ある $k,A,g$ について $h\preceq\lambda_{k,g}$ となる有限型付け $h$ を
\emph{局所群決定的}という．
\end{definition}

定義~\ref{def:prefix-suffix-typing} で $\ell=k$ とすると，
長さ $2k$ 未満の語は\emph{短い}語としてそのまま記録され，
\emph{長い}語 $w$ は $\langle\operatorname{pref}_k(w),\operatorname{suf}_k(w)\rangle$ として記録される．
$s,p\in\Sigma^k$ に対し
$J_g(s,p):=\prod_{i=1}^{k}g\bigl((sp)[i..i+k]\bigr)$ と書く．
これは，$sp$ の長さ $k+1$ の因子のうち $s$ の中で始まる $k$ 個についての順序つき積である．

\begin{lemma}[局所群型付けは準同型である]
\label{lem:local-group-morphism}
$\lambda_{k,g}$ の核は $\Sigma^*$ 上の有限指数の合同である．
したがって $\lambda_{k,g}$ は有限モノイド $M_{k,g}$ への全射準同型である．
長い $u,v$ に対して
\[
 \lambda_{k,g}(uv)
 =\bigl(\langle\operatorname{pref}_k u,\operatorname{suf}_k v\rangle,\;
   \gamma_{k,g}(u)\,J_g(\operatorname{suf}_k u,\operatorname{pref}_k v)\,
   \gamma_{k,g}(v)\bigr)
\]
が成り立つ．
\end{lemma}

\begin{proof}
窓成分が準同型であることは定義~\ref{def:prefix-suffix-typing} による．
$u,v$ が長いとき，$uv$ の長さ $k+1$ の因子は順に，
$u$ の因子，$u$ の最後の $k$ 位置で始まる $k$ 個の因子，$v$ の因子である．
中央の因子は $\operatorname{suf}_k(u)\operatorname{pref}_k(v)$ に含まれるので，表示の式を得る．

$u$ が短く $v$ が長いとき，$uv$ の因子のうち $u$ で始まるものは
$u\operatorname{pref}_k(v)$ に含まれ，残りは $v$ の因子である．
したがって $\gamma_{k,g}(uv)$ は $u$，$\operatorname{pref}_k(v)$，$\gamma_{k,g}(v)$ で決まる．
$u$ が長く $v$ が短い場合も対称である．
両方が短ければ $uv$ そのものが決まる．

いずれの場合も $\lambda_{k,g}(uv)$ は $\lambda_{k,g}(u)$ と $\lambda_{k,g}(v)$ の関数である．
像の大きさは高々 $|\Sigma|^{<2k}+|\Sigma|^{2k}|A|$ である．
\end{proof}

\begin{lemma}[ブロック分解]
\label{lem:local-group-blocks}
語を，指定されたブロック $\beta_j$ と固定語 $x_j$ により
$w=x_0\beta_1x_1\cdots\beta_qx_q$ と分解する．
各 $j$ について $h^{k,k}_\Sigma(\beta'_j)=h^{k,k}_\Sigma(\beta_j)$ となる $\beta'_j$ を用いて
$w'=x_0\beta'_1x_1\cdots\beta'_qx_q$ と置く．
このとき $h^{k,k}_\Sigma(w')=h^{k,k}_\Sigma(w)$ である．
さらに $w$ が長ければ，$x_j$ と $\beta_j$ の窓型だけに依存する
$c_0,\ldots,c_q\in A$ が存在して
\[
 \gamma_{k,g}(w)=c_0\,\gamma_{k,g}(\beta_1)\,c_1\cdots
   \gamma_{k,g}(\beta_q)\,c_q,
 \qquad
 \gamma_{k,g}(w')=c_0\,\gamma_{k,g}(\beta'_1)\,c_1\cdots
   \gamma_{k,g}(\beta'_q)\,c_q
\]
が成り立つ．
\end{lemma}

\begin{proof}
短いブロックはそのまま記録されるので $\beta'_j=\beta_j$ であり，固定語に吸収できる．
残るブロックはすべて長く，その置換も長く，最初と最後の $k$ 文字が等しい．
よって $w'$ が長いのは $w$ が長いときに限り，最初と最後の $k$ 文字は変わらない．
これで窓についての主張が従う．

$w$ の長さ $k+1$ の因子 $F$ を分類する．
$F$ が長いブロック $\beta_j$ の内部にあれば，$F$ は $\beta_j$ の因子の一つである．
そうでなければ，$F$ が接する各長いブロックとの交わりは，長さ高々 $k$ の接頭辞か接尾辞である．
$F$ が $\beta_j$ の両端に接するなら $|\beta_j|\le k-1$ となり，$\beta_j$ は短い．
よって第二種の因子はすべて，固定語と窓
$\operatorname{pref}_k(\beta_j)$，$\operatorname{suf}_k(\beta_j)$ だけで綴られ，
同じ因子が同じ順序で $w'$ にも現れる．
因子を開始位置の順に並べると，$\beta_j$ 内部の因子は連続する一続きをなし，
$\gamma_{k,g}(\beta_j)$ を与える．
$\beta_j$ と $\beta_{j+1}$ の一続きの間にある残りの因子の順序つき積を $c_j$ とする．
\end{proof}

\begin{theorem}[局所群正像の特徴づけ]
\label{thm:lg-characterization}
有限型付け $h:\Sigma^*\to M$ について $S:=h(\Sigma^+)$ と置く．次は同値である．
\begin{enumerate}[label=\textup{(\roman*)},leftmargin=*]
\item $S\in\mathbf{LG}$．すなわち，$S$ の各冪等元 $e$ について $eSe$ は群である．
\item $S$ は $U_1$ と同型な部分半群をもたない．
\item $S$ の冪等元はすべて $S$ の最小イデアルに属する．
\item $h$ は局所群決定的である．$k=|S|$ と取れる．
\end{enumerate}
\end{theorem}

\begin{proof}
(i)$\Rightarrow$(ii)．$ef=fe=f\ne e$ となる $\{e,f\}\cong U_1$ があるとする．
すると $f\in eSe$ は単位元 $e$ をもつ群の冪等元なので $f=e$ となり，矛盾する．

(ii)$\Rightarrow$(iii)．$S$ の最小イデアル $K$ は冪等元を含む．
$S$ の冪等元 $e$ について $s\in K$ を選び，$x:=(ese)^\omega$ と置く．
これは $K\cap eSe$ の冪等元なので $ex=xe=x$ である．
$x\ne e$ なら $\{e,x\}\cong U_1$ となるので，$e=x\in K$ である．

(iii)$\Rightarrow$(i)．有限半群の最小イデアル $K$ は完全単純である~\cite{Howie1995}．
冪等元 $e\in K$ について $eSe\subseteq K$ なので $eSe=eKe$ である．
完全単純半群ではこれは $e$ の $\mathcal H$-類であり，群である．

(iii)$\Rightarrow$(iv)．Rees の定理~\cite{Howie1995} により
$K=\mathcal M(H;I,\Lambda;P)$ と書く．積は
$(i,a,\lambda)(j,b,\mu)=(i,ap_{\lambda j}b,\mu)$ である．$n:=|S|$ と置く．

\emph{長い積は $K$ に入る．}
$S$ の冪等元はすべて $K$ に属するので，Rees 商 $S/K$ は $0$ 以外の冪等元をもたない．
$m:=|S/K|\le n$ として $s_1,\ldots,s_m\in S$ を取り，$S/K$ における接頭積を考える．
どれも $0$ でなければ二つが一致し，$i<j$ について $p_i=p_j$ となる．
このとき $x=s_{i+1}\cdots s_j$ について $p_i=p_ix=p_ix^\omega=0$ となり，矛盾する．
よって $|w|\ge n$ なら $h(w)\in K$ である．

\emph{右作用．}
$i_0\in I$ を固定し，$\lambda\in\Lambda$ について冪等元
$e_\lambda:=(i_0,p_{\lambda i_0}^{-1},\lambda)$ を取る．
$x=(i,a,\lambda)\in K$ について $xe_\lambda=x$ なので，
$s\in S$ に対し $xs=x(e_\lambda s)$ である．
$e_\lambda s=(j,c,\mu)\in K$ と書くと
\[
 (i,a,\lambda)\,s=(i,\;a\,\rho(\lambda,s),\;\mu(\lambda,s)),
 \qquad \rho(\lambda,s):=p_{\lambda j}c
\]
を得る．ここで $\rho(\lambda,s)\in H$ と $\mu(\lambda,s)\in\Lambda$ は $(\lambda,s)$ だけで決まる．
双対的に，$|w|\ge n$ について $h(w)$ の $\Lambda$ 座標は $h(\operatorname{suf}_n w)$ の $\Lambda$ 座標に等しく，
$I$ 座標は $h(\operatorname{pref}_n w)$ の $I$ 座標に等しい．

\emph{窓による分解．}
$m\ge2n$ について $w=a_1\cdots a_m$ とし，
$j\ge n$ について $h(a_{j-n+1}\cdots a_j)$ の $\Lambda$ 座標を $\lambda_j$ とする．
右作用を繰り返すと
\[
 h(w)=\Bigl(i,\;a\prod_{j=n}^{m-1}\rho(\lambda_j,h(a_{j+1})),\;
 \lambda_m\Bigr),
 \qquad h(a_1\cdots a_n)=(i,a,\lambda_n)
\]
となる．積の $j$ 番目の因子は，長さ $n+1$ の因子 $a_{j-n+1}\cdots a_{j+1}$ だけで決まる．
$F\in\Sigma^{n+1}$ について
$g(F):=\rho\bigl(\lambda(h(F[1..n])),h(F[n+1])\bigr)$ と定めると，この積は $\gamma_{n,g}(w)$ である．
よって $h(w)$ は $\operatorname{pref}_n w$，$\operatorname{suf}_n w$，$\gamma_{n,g}(w)$ で決まる．
長さ $2n$ 未満の語は $h^{n,n}_\Sigma$ がそのまま記録し，
$\eps$ はその $\lambda_{n,g}$-類のただ一つの語である．
以上より $A=H$ として $h\preceq\lambda_{n,g}$ である．

(iv)$\Rightarrow$(i)．$\mathbf{LG}$ は準同型像で閉じているので，
$T:=\lambda_{k,g}(\Sigma^+)\in\mathbf{LG}$ を示せばよい．
短い元は冪等でない．短い語の平方は，より長い短い語になるか，長い語になるからである．
冪等元 $e=(\langle p,s\rangle,\gamma_e)$ を取る．
補題~\ref{lem:local-group-morphism} により $\gamma_eJ_g(s,p)\gamma_e=\gamma_e$ である．
$e$ は両側で長いので，各 $x\in T$ について $exe$ の窓は $\langle p,s\rangle$ である．
よって $eTe$ は $H_e:=\{(\langle p,s\rangle,\alpha):\alpha\in A\}$ に含まれる．
$H_e$ 上の積は $\alpha*\beta=\alpha J_g(s,p)\beta$ であり，
$\alpha\mapsto\alpha J_g(s,p)$ は $(H_e,*)$ から $A$ への同型である．
したがって $eTe$ は有限群の部分半群であり，群である．
\end{proof}

\begin{proposition}[窓・群型付けは局所群型付けであり，包含は真である]
\label{prop:li-g-strict}
\begin{enumerate}[label=\textup{(\roman*)},leftmargin=*]
\item 有限群への準同型 $g$ について $h\preceq(h^{k,\ell}_\Sigma,g)$ なら，
$K=\max\{k,\ell,1\}$，$g'(F):=g(F\text{ の最後の文字})$ として $h\preceq\lambda_{K,g'}$ である．
したがって $\mathbf{LI}\vee\mathbf G\subseteq\mathbf{LG}$ である．
\item $\Sigma=\{\mathtt a,\mathtt b,\mathtt c\}$ とし，
$g(\mathtt{ab})=1\in\mathbb Z_2$，それ以外の $F$ について $g(F)=0$ として
$h_{\mathtt{ab}}:=\lambda_{1,g}$ と置く．
これは窓とともに $\mathtt{ab}$ の出現回数の偶奇を記録する．
$h_{\mathtt{ab}}$ は窓・群決定的でない．
よって $h_{\mathtt{ab}}(\Sigma^+)\in\mathbf{LG}\setminus(\mathbf{LI}\vee\mathbf G)$ である．
\end{enumerate}
\end{proposition}

\begin{proof}
(i) $|w|\ge2K$ なら $g(w)=g(\operatorname{pref}_Kw)\,\gamma_{K,g'}(w)$ である．
それより短い語はそのまま記録される．
pseudovariety の包含は，定理~\ref{thm:lg-characterization} と，
$\mathbf{LI}\vee\mathbf G$ を窓・群型付けで特徴づける姉妹論文の結果~\cite[Lemma~7.4]{kuriyamaCFG} から従う．

(ii) $g':\Sigma^*\to A'$ について $h_{\mathtt{ab}}\preceq(h^{k,\ell}_\Sigma,g')$ と仮定する．
$K=\max\{k,\ell,1\}$ とし，$q$ を $A'$ の冪数の偶数倍とする．
$u=\mathtt a^K(\mathtt{abc})^q\mathtt a^K$，
$v=\mathtt a^K(\mathtt{bac})^q\mathtt a^K$ と置く．
どちらも長く，接頭辞 $\mathtt a^k$，接尾辞 $\mathtt a^\ell$ をもち，
$g'(u)=g'(\mathtt a)^{2K}=g'(v)$ である．
しかし $u$ は $\mathtt{ab}$ を $q$ 回含むのに対し，
$v$ はちょうど1回（$\mathtt a^K\mathtt b$ の位置）しか含まない．
よって $h_{\mathtt{ab}}(u)\ne h_{\mathtt{ab}}(v)$ となり，矛盾する．
$\mathbf{LG}$ への所属は定理~\ref{thm:lg-characterization} から，
$\mathbf{LI}\vee\mathbf G$ への非所属は姉妹論文の特徴づけから従う．
\end{proof}

\begin{lemma}[局所群型を指定した短いタプル]
\label{lem:local-group-shortest}
$G=(V,\Sigma,P,S,\fan)$ を，ファンアウト高々 $f$，規則ランク高々 $r_{\max}\ge1$ の
縮約済み良形 MCFG とし，$\lambda=\lambda_{k,g}$ とする．
縮約された $h^{k,k}_\Sigma$-精密化の開始記号以外の型付き非終端記号の個数を
$n_{\mathrm w}\le|V||M_{k,k}|^f$ とし，
$b_{\mathrm w}:=f(2k+1)n_{\mathrm w}(r_{\max}+1)\bar\theta_G$ と置く．
$n_\lambda:=|V||M_{k,g}|^f$ とし，$M_{A,f}$ を補題~\ref{lem:group-heavy-antichain} の定数とする．
各 $B\in V$ と実現可能な各 $\mathbf t\in M_{k,g}^{\fan(B)}$ について，
$\lambda^{(\fan(B))}(\tuple u)=\mathbf t$ かつ
\[
 \|\tuple u\|_1\le\beta_\lambda:=(r_{\max}+1)\,b_{\mathrm w}
 \max\{1,(M_{A,f}-1)\,n_\lambda\}
\]
となる $\tuple u\in L_B(G)$ が存在する．
$A$ が可換なら $M_{A,f}-1$ を $\mathsf d(A^f)$ に置き換えてよい．
\end{lemma}

\begin{proof}
根の型が $\mathbf t$ となる $B$ からの終端導出のうち，
まず総出力長，次に節点数を最小化するもの $T$ を選ぶ．
各節点に，非終端記号とその出力の $\lambda$-ベクトルを標識する．
補題~\ref{lem:group-typed-shortest} と同様に，経路上で標識は繰り返されない．
$\lambda$ は準同型であり（補題~\ref{lem:local-group-morphism}），
規則テンプレートは準同型的に評価されるからである．
よって各経路の節点数は高々 $n_\lambda$ である．

出力の総長が $b_{\mathrm w}$ を超える節点を\emph{重い}と呼ぶ．重い節点は祖先について閉じている．
重い節点 $v$ の標識を $C$ とし，その出力の $h^{k,k}_\Sigma$-ベクトルを $\boldsymbol\omega$ とする．
$G$ は縮約済みなので，$T$ は成功する導出に埋め込める．
よって $C_{\boldsymbol\omega}$ は縮約された $h^{k,k}_\Sigma$-精密化に残る型付き非終端記号であり，
補題~\ref{lem:boundary-short-canonical} により，
出力の窓ベクトルが同じで総長が高々 $b_{\mathrm w}$ の $C$ からの導出が存在する．
各重い節点についてこの置換を固定する．

重い節点の反鎖 $v_1,\ldots,v_m$ を考える．
各 $v_i$ は根の成分に高々 $f$ 個の互いに素なブロックを与える．
任意の添字集合 $J$ の部分木を置換すると，総長が真に小さい正しい導出が得られる．
補題~\ref{lem:local-group-blocks} を根の各成分に適用すると，
この置換は根の窓型を変えない．
また長い根成分の $\gamma$ 値は，各ブロックが固定定数の間に自分の $\gamma$ 値を与える
順序つき積を通じてのみ変わる．
これを群$A$上のブロック置換族として明示する．
各添字$i$は節点$v_i$が根の各成分に与えるブロックを所有し，
その総数は高々$f$である．窓型の短いブロックは
型一致により実際の文字列が不変だから，固定部分へ吸収する．
残る長いブロックの$\gamma$-値だけが置換によって変化し，
補題~\ref{lem:local-group-blocks}の定数因子は選択集合$J$に依存しない．
したがって各根成分の$\gamma$-値は，
固定定数と各所有者の元のブロック値／置換後のブロック値との
非可換な\emph{順序付き積}として表せる．
所有するブロックがゼロの添字があればそれだけを置換でき，
根の型を保って出力を短縮できるので，
以下では各添字が少なくとも一つのブロックを持つと仮定してよい．

所有者ごとのブロック数，各ブロックの根成分，
同一所有者内の順序で定まるパターンは高々
$U_f=f^{f+1}f!$種類である．
$m\ge M_{A,f}=U_fN_{A,f}^{2^{f^2}}$なら，
同一パターンの添字が少なくとも
$N_{A,f}^{2^{f^2}}$個存在する．
その他のブロックは固定定数へ吸収する．
補題~\ref{lem:order-homogenize}を高々$f^2$個の
比較型へ適用すると，少なくとも
$N_{A,f}=\max\{3,|A|^f\}$個の添字について
所有者間のブロックの順序比較をすべて斉次化できる．
その部分族には補題~\ref{lem:group-interval}が適用でき，
非空の添字区間を同時に置換して各根成分の
順序付き群積を完全に保存する．
窓型も保存されるため根の$\lambda$-ベクトルが一致し，
各重い節点の置換は総長を真に減らすので最小性に矛盾する．
$m\ge M_{A,f}$ なら，置換のある非空集合が根の $\lambda$-ベクトル全体を保存し，最小性に矛盾する．
$A$ が可換なら，補題~\ref{lem:group-typed-shortest} の零和論法を差分
$\gamma(\beta')-\gamma(\beta)$ に適用して，上界 $\mathsf d(A^f)$ を得る．

以上により，任意の重い節点反鎖の濃度は
高々$M_{A,f}-1$である．
この条件を，$\lambda=\lambda_{k,g}$，
$b=b_{\mathrm w}$，$c=M_{A,f}-1$として
補題~\ref{lem:heavy-node-count}に適用すれば，
根タプルの総長は表示した$\beta_\lambda$以下となる．
可換群の場合はブロック置換の零和論法により
$c=\mathsf d(A^f)$と置ける．
\end{proof}

\begin{theorem}[局所群型付けに対する多項式データ]
\label{thm:local-group-polydata}
$f,r_{\max},\Sigma$ と，$h(\Sigma^+)\in\mathbf{LG}$ となる有限型付け $h$ を固定する．
任意の $L\in\Cmcf{f}{h}$ と，$L(G)=L\cap\Sigma^+$ を満たす
ファンアウト高々 $f$，ランク高々 $r_{\max}$ の縮約済み良形文法 $G$ について，上限つき作用素 $\mathcal B^{\le r_{\max}}_{f,h}$ は
\[
 |\Wit(\widetilde G_0)|=O(\|G\|),\qquad
 \|\Wit(\widetilde G_0)\|_+=O_{f,r_{\max},h,\Sigma}(\|G\|^3\bar\theta_G)
\]
となる特徴標本をもつ．
\end{theorem}

\begin{proof}
定理~\ref{thm:lg-characterization} により，$h$ だけに依存するある $k,A,g$ について
$h\preceq\lambda_{k,g}$ である．
$X=B_{\mathbf s}$ を縮約された $h$-精密化の開始記号以外の型付き非終端記号とし，
$\tuple u\in L_X$ とする．
$\tuple u$ の $\lambda$-ベクトルに補題~\ref{lem:local-group-shortest} を適用すると，
同じ $\lambda$-ベクトル（したがって同じ $h$-ベクトル）をもち
$\|\tuple u'\|_1\le\beta_\lambda$ となる $\tuple u'\in L_B(G)$ が得られる．
命題~\ref{prop:output-invariants}(ii) によりその導出は根 $X$ に持ち上がり，
補題~\ref{lem:boundary-short-canonical} と同様に縮約後も残る．
よって長さ高々 $\beta_\lambda=O(\|G\|^2\bar\theta_G)$ の正準タプルが存在する．
ここで $b_{\mathrm w}=O(\|G\|\bar\theta_G)$，$n_\lambda=O(\|G\|)$ を用いた．
補題~\ref{lem:untyped-contexts} により各証拠語の長さは高々
$2\bar\theta_G+r_{\max}\beta_\lambda$ であり，
命題~\ref{prop:witness-cardinality} によりその個数は $O(\|G\|)$ である．
\end{proof}

定理~\ref{thm:local-group-polydata} は，窓型付けと窓・群型付けの定理を特別な場合として含み
（命題~\ref{prop:li-g-strict}(i)），群の場合と同じ指数3を与える．
下界側は次の交わり補題によって正確に述べられる．

\begin{lemma}[型フィルタは代入可能性を保つ]
\label{lem:type-filter}
$L$ が $(f,h)$-タプル代入可能で $F\subseteq M$ なら，
$L\cap h^{-1}(F)$ も $(f,h)$-タプル代入可能である．
\end{lemma}

\begin{proof}
$\tuple x,\tuple y$ が等しい $h$-ベクトルをもち，
$L':=L\cap h^{-1}(F)$ において文脈 $E$ を共有するとする．
両者は $L$ でも $E$ を共有するので，$L$ における分布が等しい．
$E'[\tuple x]\in L'$ なら $E'[\tuple y]\in L$ であり，
$h(E'[\tuple y])=h(E'[\tuple x])\in F$ なので $E'[\tuple y]\in L'$ である．
\end{proof}

\begin{corollary}[pseudovariety の二分法]
\label{cor:pseudovariety-dichotomy}
$\mathbf V$ を有限半群の pseudovariety とする．
\begin{enumerate}[label=\textup{(\roman*)},leftmargin=*]
\item $\mathbf V\subseteq\mathbf{LG}$ なら，任意のアルファベット，
$h(\Sigma^+)\in\mathbf V$ となる任意の型付け $h$，任意の固定された $f,r_{\max}$ について，
定理~\ref{thm:local-group-polydata} のとおり総長 $O(\|G\|^3\bar\theta_G)$ の特徴データが存在する．
\item $\mathbf V\not\subseteq\mathbf{LG}$ なら $U_1\in\mathbf V$ である．
ファンアウト1において，$h_{\mathrm U}(\Sigma^+)=U_1$ となる型付け $h_{\mathrm U}$ と，
$\|G\|,\theta_G=O(n)$ の標的が存在し，その標的の
$\mathcal B^{\le2}_{1,h_{\mathrm U}}$ に対する特徴標本はどれも長さ $3\cdot2^n-2$ の語を含む．
$\Cmcf{1}{h_{\mathrm U}}$ 全体で特徴標本をもつ任意の集合駆動作用素についても同じである．
\end{enumerate}
\end{corollary}

\begin{proof}
(i) は定理~\ref{thm:local-group-polydata} である．
(ii) について，$S\in\mathbf V\setminus\mathbf{LG}$ を選ぶ．
定理~\ref{thm:lg-characterization} により $S$ は $U_1$ と同型な部分半群をもち，
pseudovariety は部分半群で閉じている．

\cite[Proposition~6.4]{kuriyamaCFG} の倍化族 $L_n$，型付け $h_{\mathrm U}$，
型1の長い語 $w_n^*$ を取る．
そこでは $L_n$ が Clark--Eyraud 代入可能であることが示されており，
したがって $(1,h_{\mathrm U})$-タプル代入可能である．
補題~\ref{lem:type-filter} により
$L_n^-:=L_n\cap h_{\mathrm U}^{-1}(\zeta)=L_n\setminus\{w_n^*\}$ もそうである．
$K\subseteq L_n$ が $w_n^*$ を含まなければ $K\subseteq L_n^-$ であり，
命題~\ref{prop:reconstruction-soundness} により
$L(\mathcal B^{\le2}_{1,h_{\mathrm U}}(K))\subseteq L_n^-$ となる．
よってどの特徴標本も $w_n^*$ を含む．
集合駆動作用素についての主張は，対 $L_n,L_n^-$ に対する
\cite[Proposition~6.4(iv)]{kuriyamaCFG} の区別対論法である．
\end{proof}

\begin{remark}[二分法の適用範囲]
\label{rem:dichotomy-scope}
\begin{enumerate}[label=\textup{(\roman*)},leftmargin=*]
\item (i) はすべての固定ファンアウトで成り立つが，(ii) はファンアウト1でしか示されていない．
倍化族はそのままでは持ち上がらない．
実際 $L_2$ は $(2,h_{\mathrm U})$-タプル代入可能でない．
型1の対 $(\mathtt[,\mathtt a\mathtt])$ と $(\mathtt[\mathtt a,\mathtt])$ は
文脈 $\mathtt{[[aa]}\,\blank_1\,\mathtt a\,\blank_2\,\mathtt]$ を共有するが，分布が異なる．
これは $n\le3$ について網羅的に確かめた．
各成分の文字集合を記録する，より細かい型付け（像はやはり半束）でも，$n=3$ で同様の破綻が現れる．
(ii) がファンアウト $f\ge2$ でも成り立つかは，
未解決問題~\ref{op:fanout-two} の pseudovariety 版である．
\item この二分法は pseudovariety についてのものである．
$h(\Sigma^+)\notin\mathbf{LG}$ となる個々の型付け $h$ について，
$\Cmcf{1}{h}$ が倍化族の符号化された複製を含むとは限らない．
そのような型付けがすべて超多項式の特徴データを強いるかは未解決である．
\end{enumerate}
\end{remark}
